\section{Domination with a smaller background}

Throughout this section $\theta$ is semipositive and big. Let $0<\delta<1$. When a singular potential lies above a $\delta\theta$-plurisubharmonic function, the gap between them can be measured by an envelope in the remaining background $(1-\delta)\theta$. This background is again semipositive and big, so its bounded potentials have Monge--Amp\`ere measures of positive mass by \cref{prop:big-noncollapsing}, and bounded domination applies to them. It bounds the envelope from below. A family of stretched envelopes then converts a strict measure inequality into a pointwise inequality, as in the bounded case of \cref{thm:semipositive-bounded-domination}.

We say that $u\in\PSH(X,\theta)$ \emph{admits a smaller-background envelope} if $P_{\delta\theta}(u)\not\equiv-\infty$ for some $0<\delta<1$. This property depends only on the ordinary singularity type of $u$: if $u'\geq u-C$, then $P_{\delta\theta}(u)-C$ is a competitor for $P_{\delta\theta}(u')$.

\Needspace{16\baselineskip}
\begin{lemma}\label{lem:singular-zero-domination}
Suppose $u,v\in\PSH(X,\theta)$, $v\preceq u$, and
$\mathrm{MA}^{\mathrm{ext}}_\theta(u)(\{u<v\})=0$. If $0<\delta<1$ and $v_\delta=P_{\delta\theta}(v)\not\equiv-\infty$, then $u\geq v_\delta$.
\end{lemma}
\begin{proof}
Choose $C$ with $v-C\leq u$. Since $v_\delta\leq v$, the obstacle $u-v_\delta$ is bounded below by $-C$ off the pole sets, and it is finite almost everywhere. Hence the envelope
\[
 p=P_{(1-\delta)\theta}(u-v_\delta)
\]
is bounded by \cref{lem:bounded-below-envelope}, and by contact support (\cref{thm:envelope-contact}) its measure is carried by $D=\{p+v_\delta=u\}$. Moreover $p+v_\delta\leq u$ everywhere, and
\[
 \theta+\ddc(p+v_\delta)
 =((1-\delta)\theta+\ddc p)+(\delta\theta+\ddc v_\delta)
\]
is a sum of two positive currents. Positivity of mixed products and ordered contact therefore give
\[
 \mathbf1_D\mathrm{MA}_{(1-\delta)\theta}(p)
 \leq\mathbf1_D\mathrm{MA}^{\mathrm{ext}}_\theta(u);
\]
this inequality is first proved on common finite truncation levels and then extended by exhaustion. Off the pole sets, every point of $D\cap\{p<0\}$ satisfies
$u=p+v_\delta<v_\delta\leq v$. The hypothesis therefore gives
$\mathrm{MA}_{(1-\delta)\theta}(p)(\{p<0\})=0$. Since $(1-\delta)\theta$ is semipositive and big, bounded domination (\cref{thm:semipositive-bounded-domination}, with \cref{prop:big-noncollapsing}), applied to $p$ and the constant $0$ with coefficient zero, yields $p\geq0$. Hence $u\geq p+v_\delta\geq v_\delta$.
\end{proof}

\begin{theorem}\label{thm:singular-strict-domination}\index{Domination principle!singular potentials}
Let $u\in\PSH(X,\theta)$ satisfy $P_{\delta\theta}(u)\not\equiv-\infty$ for some $0<\delta<1$. Suppose $v\in\PSH(X,\theta)$, $v\preceq u$, and $0\leq c<1$. If
\[
 \mathbf1_{\{u<v\}}\mathrm{MA}^{\mathrm{ext}}_\theta(u)
 \leq c\,\mathbf1_{\{u<v\}}\mathrm{MA}^{\mathrm{ext}}_\theta(v),
\]
then $u\geq v$.
\end{theorem}
\begin{proof}
Since $\max(u,v)=v$ on the plurifine open set $\{u<v\}$, plurifine locality shows that replacing $v$ by $\max(u,v)$ preserves the hypothesis; it does not change the conclusion. We may therefore assume $u\leq v$, and since $v\preceq u$ there is $C\geq0$ with
$u\leq v\leq u+C$. For $b>1$ consider the stretched envelopes
\[
 p_b=P_\theta(bu-(b-1)v),\qquad
 D_b=\{p_b=bu-(b-1)v\}.
\]
The function $u-(b-1)C$ is a competitor, so $p_b\not\equiv-\infty$; and since $bu-(b-1)v\leq u$, we have
$u-(b-1)C\leq p_b\leq u$. By contact support, the measure of $p_b$ is concentrated on $D_b$. The potential
$b^{-1}p_b+(1-b^{-1})v$ lies below $u$ and agrees with it on $D_b$, so ordered contact and positivity of mixed products give
\[
 \mathbf1_{D_b}\bigl[b^{-n}\mathrm{MA}^{\mathrm{ext}}_\theta(p_b)
 +(1-b^{-1})^n\mathrm{MA}^{\mathrm{ext}}_\theta(v)\bigr]
 \leq\mathbf1_{D_b}\mathrm{MA}^{\mathrm{ext}}_\theta(u).
\]
Choose $b$ so large that $(1-b^{-1})^n>c$. On $D_b\cap\{u<v\}$ the hypothesis then forces $\mathrm{MA}^{\mathrm{ext}}_\theta(p_b)$ to vanish. Since the masses may be infinite, we justify this subtraction as follows. On each set where $u,v,p_b>-k$, every measure involved is finite, being bounded by the mass of the corresponding bounded truncation. We subtract the finite terms on each such set and then take the union over $k$, so no infinite masses are ever subtracted.

It follows that the measure of $p_b$ is carried by $D_b\cap\{u=v\}\subset\{p_b=u\}$. We apply \cref{lem:singular-zero-domination} to the pair $(p_b,u)$: their singularities are equivalent, $\mathrm{MA}^{\mathrm{ext}}_\theta(p_b)(\{p_b<u\})=0$, and $P_{\delta\theta}(u)\not\equiv-\infty$. We obtain
\[
 p_b\geq P_{\delta\theta}(u)
\]
for all sufficiently large $b$. Suppose now that $u\geq v$ fails. By \cref{lem:qpsh-almost-everywhere-order}, the set $\{u<v\}$ has positive volume. Since both potentials are finite almost everywhere, some set $E=\{u<v-a\}$ with $a>0$ has positive volume and both functions are finite on $E$. On $E$ the obstacle inequality gives
\[
 P_{\delta\theta}(u)\leq p_b\leq v-ba
\]
quasi-everywhere. Letting $b$ tend to infinity through the integers shows that the quasi-plurisubharmonic function $P_{\delta\theta}(u)\not\equiv-\infty$ equals $-\infty$ on a set of positive volume, a contradiction. Thus $u\geq v$ off the pole sets, and \cref{lem:qpsh-almost-everywhere-order} gives the inequality everywhere.
\end{proof}

\begin{corollary}\label{cor:singular-exponential-order}
Let $u,v\in\PSH(X,\theta)$ satisfy $v\preceq u$, and suppose that $u$ admits a smaller-background envelope. If $\lambda>0$ and
\[
 \mathrm{e}^{-\lambda u}\mathrm{MA}^{\mathrm{ext}}_\theta(u)
 \leq \mathrm{e}^{-\lambda v}\mathrm{MA}^{\mathrm{ext}}_\theta(v)
\]
off the pole sets, then $u\geq v$. In particular, for every $\lambda>0$ and every positive measure $\mu$, the equation $\mathrm{MA}_\theta(u)=\mathrm{e}^{\lambda u}\mu$ has at most one solution in an ordinary singularity type admitting a smaller-background envelope.
\end{corollary}
\begin{proof}
Fix $t>0$. On $\{u<v-t\}$ the assumed inequality gives
$\mathrm{MA}_\theta^{\mathrm{ext}}(u)\leq \mathrm{e}^{-\lambda t}\mathrm{MA}_\theta^{\mathrm{ext}}(v-t)$.
The theorem, applied to $u$ and $v-t$ with $c=\mathrm{e}^{-\lambda t}$, gives $u\geq v-t$; now let $t\downarrow0$. For two solutions of the same equation in the stated type, apply this result in both directions. The weighted measures are compared off the common pole sets, which these measures do not charge, so no value of an infinite weight is needed.
\end{proof}
